\chapter{AKSZ--RSVP Action and Classical Master Equation}

\section{Overview and Strategy}

The goal of this chapter is to apply the AKSZ construction of
Chapter~4 to the derived lamphrodynamic plenum constructed in
Chapter~6. The result is a $(-1)$-shifted symplectic BV theory whose
classical Euler--Lagrange equations agree with the lamphrodynamic
equations of Chapter~5. This constitutes the precise mathematical
definition of the RSVP field theory at the classical level.

Let $(M,g)$ be a smooth oriented time-oriented $4$-dimensional Lorentzian
manifold with associated shifted tangent bundle $\mathsf{T}[1]M$.
Define the AKSZ field space
\[
\mathcal{F} := \mathrm{Map}(\mathsf{T}[1]M,\mathrm{Plenum}),
\]
a derived mapping stack equipped with the canonical $(-1)$-shifted
symplectic form of Chapter~4.
Fields in $\mathcal{F}$ consist of superfields whose degree-zero
components are the lamphrodynamic fields $(\Phi,\vec v,S)$ and whose
higher-degree components constitute ghosts, antifields, and higher
ghosts.


\section{AKSZ Data for RSVP}

The AKSZ construction requires:

\begin{enumerate}
\item a source $(\mathsf{T}[1]M,d_M)$ with homological vector field
$d_M$, the de~Rham differential on $M$ shifted by $1$;
\item a target $(\mathrm{Plenum},\omega)$ with $2$-shifted symplectic
form $\omega$ and homological vector field $Q_{\mathrm{Plenum}}$
preserving $\omega$;
\item a Hamiltonian function $\alpha$ on $\mathrm{Plenum}$ satisfying
$Q_{\mathrm{Plenum}}=\{\alpha,-\}$.
\end{enumerate}

Under these assumptions, the AKSZ recipe produces the BV action
\[
S_{\mathrm{AKSZ}} = \int_{\mathsf{T}[1]M}
\langle \Phi^{\ast}\theta, d_M\Phi\rangle + \Phi^{\ast}\alpha,
\]
where $\theta$ is a potential for the shifted symplectic form on
$\mathrm{Plenum}$.
Explicit formulas for $\theta$ will be given below (Section~\ref{sec:theta}).


\section{Derived 3-Form and BV Potential}

Recall from Chapter~6 the derived $3$-form $\Theta$ on $\mathrm{Plenum}$
given locally by
\[
\Theta =
\Phi\, dS\wedge \star d\Phi
+ \beta\, \langle\nabla\Phi,\nabla S\rangle\, dV_g
+ \gamma\, \langle\nabla S,\vec v\rangle\, d(\iota_{\vec v}dV_g).
\]
Let $\theta$ be any degree-$2$ potential for $\omega=d_{\mathrm{dR}}\Theta$
in the derived de~Rham complex of $\mathrm{Plenum}$. Such a potential
exists locally and may be constructed globally by descent, using the
Artin atlases of $\mathrm{Plenum}$.

\begin{remark}
In practice, $\theta$ may be chosen to be the $2$-form part of $\Theta$,
omitting terms that vanish under pullback to the mapping stack.
However, the existence of a global potential is essential for defining
the AKSZ action.
\end{remark}


\section{The RSVP AKSZ Action}
\label{sec:theta}

The AKSZ action is now defined by
\begin{equation}
\label{eq:RSVP-AKSZ}
S_{\mathrm{RSVP}}
:= \int_{\mathsf{T}[1]M}
\left(
\langle \Phi^{\ast}\theta, d_M\Phi\rangle
+ \Phi^{\ast}\alpha
\right).
\end{equation}
In local coordinates, writing $\Phi^{\ast}=(\Phi,\vec v,S)$ for the
pullback of the universal fields, we may express the integrand as
\begin{align}
\mathcal{L}_{\mathrm{AKSZ}}
&=
\langle d_M\Phi, \theta_{\Phi}\rangle
+ \langle d_M \vec v, \theta_{v}\rangle
+ \langle d_M S, \theta_{S}\rangle
+ \alpha(\Phi,\vec v,S),
\end{align}
where $\theta_{\Phi},\theta_{v},\theta_{S}$ are the components of the BV
potential paired with the corresponding superfield derivatives.

\begin{proposition}
The functional $S_{\mathrm{RSVP}}$ is well defined on the derived
mapping stack $\mathcal{F}$ and has total ghost number zero.
\end{proposition}

\begin{proof}[Proof (sketch)]
Total degree zero follows from degree counting: the degree of
$d_M$ is $+1$, the degree of $\theta$ is $2$, and the integration over
$\mathsf{T}[1]M$ reduces degree by $3$, yielding degree $0$ overall.
See \cite{AKSZ97,CF07}.
\end{proof}

\section{Classical Master Equation}

Let $\{-,-\}_{\mathrm{BV}}$ denote the BV antibracket induced by the
$(-1)$–shifted symplectic structure on $\mathcal{F}
= \mathrm{Map}(\mathsf{T}[1]M,\mathrm{Plenum})$. The classical master
equation is
\[
\{S_{\mathrm{RSVP}},S_{\mathrm{RSVP}}\}_{\mathrm{BV}} = 0.
\]

\begin{theorem}[Classical Master Equation]
\label{thm:CME-RSVP}
The RSVP AKSZ action \eqref{eq:RSVP-AKSZ} satisfies the classical master
equation
\[
\{S_{\mathrm{RSVP}},S_{\mathrm{RSVP}}\}_{\mathrm{BV}} = 0,
\]
and hence defines a classical BV theory.
\end{theorem}

\begin{proof}[Proof (Sketch)]
This follows from Theorem~\ref{thm:Main-Plenum}
and the general AKSZ theorem of
Alexandrov–Kontsevich–Schwarz–Zaboronsky:
whenever the target carries a $k$–shifted symplectic structure with
Hamiltonian $Q_{\mathrm{Plenum}}=\{\alpha,-\}$,
the induced action on
$\mathrm{Map}(\mathsf{T}[1]M,\mathrm{Plenum})$
satisfies the classical master equation and
induces the correct BV differential $Q=\{S,-\}$.
\end{proof}

\begin{remark}
Note that no gauge-fixing is used. The ghosts and antifields arise
canonically from the graded geometry of $\mathrm{Map}(\mathsf{T}[1]M,
\mathrm{Plenum})$ and the shifted symplectic structure.
\end{remark}


\section{BV Operator and Field Content}
\label{sec:BV-operator}

The BV differential is defined by
\[
Q_{\mathrm{BV}}(F) :=
\{ S_{\mathrm{RSVP}}, F\}_{\mathrm{BV}}, \qquad F\in
\mathcal{O}(\mathcal{F}).
\]
The field content consists of the superfields
\[
\mathbf{\Phi} =
\Phi + \Phi^{+} + \Phi^{\ast} + \cdots,
\qquad
\mathbf{v}
= \vec v + v^{+} + v^{\ast} + \cdots,
\qquad
\mathbf{S}
= S + S^{+} + S^{\ast} + \cdots,
\]
with ghost number assignments determined by the degree shifts of
$\mathsf{T}[1]M$ and $\mathrm{Plenum}$. The precise assignments follow
the standard AKSZ convention:
\[
\mathrm{gh}(\mathrm{field})=0,\qquad
\mathrm{gh}(\mathrm{ghost})=+1,\qquad
\mathrm{gh}(\mathrm{antifield})=-1,\qquad
\mathrm{gh}(\mathrm{higher~ghosts})>+1.
\]

\begin{proposition}
\label{prop:BV-square-zero}
The BV operator $Q_{\mathrm{BV}}$ is a degree $+1$ derivation satisfying
$Q_{\mathrm{BV}}^2=0$.
\end{proposition}

\begin{proof}[Proof (Sketch)]
$Q_{\mathrm{BV}}^2(F)=\{S,\{S,F\}\} = \frac12\{\{S,S\},F\}$, which
vanishes by Theorem~\ref{thm:CME-RSVP}.
\end{proof}


\section{Euler–Lagrange Equations and Lamphrodynamics}

We now extract the classical physical equations of motion as the
Euler–Lagrange equations of the AKSZ action at ghost number zero.

\begin{theorem}[RSVP Equations from AKSZ]
Let $S_{\mathrm{RSVP}}$ be the AKSZ action
\eqref{eq:RSVP-AKSZ}. Then the ghost-number-zero Euler--Lagrange
equations are equivalent to the lamphrodynamic field equations of
Chapter~5: namely,
\[
\Box \Phi = \lambda\nabla\!\cdot\vec{v} + \dots,
\qquad
\nabla\cdot \vec v = f(\nabla\Phi,\nabla S),
\qquad
\partial_t S + \nabla\!\cdot(S\vec v) = \kappa\nabla^2 S + \dots,
\]
with all nonlinear source terms determined by the same derived $3$–form
$\Theta$ used in the AKSZ construction.
\end{theorem}

\begin{proof}[Proof (Sketch)]
The variation of $S_{\mathrm{RSVP}}$ with respect to
$(\Phi,\vec v,S)$ at ghost number zero yields the classical field
equations. All antifield and ghost terms vanish at ghost number zero,
leaving precisely the lamphrodynamic PDE system of Chapter~5.
\end{proof}

\begin{remark}
Thus the AKSZ construction is not merely compatible with
lamphrodynamics; it {\em determines} the lamphrodynamic equations as the
classical sector of a $(-1)$–shifted symplectic BV theory. The derived
geometry forces the PDE system rather than assuming it.
\end{remark}


\section{Physical Interpretation}

The significance of the AKSZ construction is that it
endows RSVP with a canonical BV quantization framework without requiring
gauge-fixing or ad-hoc ghost terms: the ghost sector is automatic.
Moreover, the classical lamphrodynamic dynamics are precisely the
Euler--Lagrange equations of the BV action, making RSVP a field theory
in the strongest modern sense. Finally, the $(-1)$–shift ensures that
the BV bracket is well-defined and that the master equation holds,
a prerequisite for any rigorous quantum theory built on the BV
formalism.

\section{Boundary Data and BV--BFV Structures}
\label{sec:RSVP-BFV}

On a manifold with boundary $\partial M\neq\emptyset$, the AKSZ
construction produces a canonical BV--BFV structure in the sense of
Cattaneo--Mnev--Reshetikhin. Concretely, the boundary data
consist of a BFV space of boundary fields $\mathcal{F}_{\partial}$ and
a degree $0$ \emph{BFV action} $S_{\partial}$ which satisfies
\[
Q_{\mathrm{BFV}}^2=0,\qquad
Q_{\mathrm{BFV}}(S_{\partial})=0,
\]
and whose cohomology governs physical boundary degrees of freedom.

\begin{theorem}[Canonical BV--BFV Boundary Structure]
\label{thm:RSVP-BFV}
Let $M$ be a $4$--manifold with (possibly empty) boundary.
Then the AKSZ action $S_{\mathrm{RSVP}}$
induces a canonical BV--BFV structure
$(\mathcal{F}_{\partial},S_{\partial})$ on $\partial M$.
Moreover, $S_{\partial}$ determines the physical boundary excitations of
the lamphrodynamic fields $(\Phi,\vec v,S)$, including entropic flux
across $\partial M$.
\end{theorem}

\begin{proof}[Proof (Sketch)]
Immediate from the general BV--BFV theorem for AKSZ models, using that
the target is $(-1)$--shifted symplectic and that $M$ has dimension $4$.
The boundary inherits a degree-$0$ Hamiltonian system governing the
physical boundary data.
\end{proof}

\begin{remark}[Physical Meaning]
The presence of a canonical BV--BFV boundary theory implies that
entropic flux and vector circulation admit well-defined ``boundary''
degrees of freedom even before gauge-fixing.
This suggests a direct route to black-hole and horizon physics once
appropriate boundary conditions are chosen (Volume~II).
\end{remark}



\section{Local Coordinates and Expansion}
\label{sec:local-coords}

Let $(x^{\mu})_{\mu=0}^{3}$ be local coordinates on $M$ and
let $(\Phi,\vec v,S)$ denote the degree-$0$ components of the superfields.
Expanding the AKSZ action to leading (classical) order gives
\begin{equation}
\label{eq:local-RSVP}
S_{\mathrm{RSVP}} \simeq
\int_{M} \Big(
\partial_{\mu}\Phi \, \Theta^{\mu\nu\rho}
\, \partial_{\nu} v_{\rho}
+ \partial_{\mu}S \, \Theta^{\mu\nu\rho}
\, \partial_{\nu} \Phi
+ \dots
\Big)\, d^{4}x,
\end{equation}
where $\Theta^{\mu\nu\rho}$ are the components of the closed
$3$--form representing the $(-1)$--shifted symplectic structure
(cf. Theorem~\ref{thm:Main-Plenum}).

\begin{remark}
Subleading terms encode the nonlinear couplings responsible for
entropy flow and lamphrodynamic vorticity. These terms become crucial
in the strong-flux regime and near entropic singularities (Ch.~12).
\end{remark}



\section{Comparison to Chern--Simons AKSZ}
\label{sec:CS-comparison}

A useful analogy is provided by $3$--dimensional Chern--Simons theory,
which is the prototypical AKSZ model in dimension $3$ with target
$\mathfrak{g}[1]$.

\subsection{Chern--Simons AKSZ recap}
Let $N$ be a $3$--manifold and $\mathfrak{g}$ a Lie algebra.
The AKSZ construction gives the classical CS action
\[
S_{\mathrm{CS}} = \int_{N}
\Big( A\wedge dA + \tfrac{1}{3}A\wedge [A,A] \Big),
\]
arising from the canonical symplectic structure on $\mathfrak{g}[1]$
and the Chevalley--Eilenberg differential.

\subsection{Structural analogy}
Two structural parallels are immediate:

\begin{enumerate}
\item
{\bf Degree shift.}
CS uses a $0$--shifted target $\mathfrak{g}[1]$, whereas RSVP uses a
$(-1)$--shifted derived stack.
The lower shift reflects the presence of physical scalar
and vector fields rather than a pure gauge sector.

\item
{\bf Closed form generating the action.}
In both theories, the action is generated from the geometric
data of the target:
the Lie algebra $(\mathfrak{g},[-,-])$ in CS,
the lamphrodynamic $3$--form $\Theta$ in RSVP.
\end{enumerate}

\subsection{Physical meaning}
For Chern--Simons, the AKSZ construction explains why the CS action is
automatically topological and why its BV–BFV boundary theory captures
edge modes (e.g.\ quantum Hall).
In RSVP the AKSZ construction explains why lamphrodynamic equations
arise as Euler--Lagrange equations of a $(-1)$--shifted BV theory,
and thus why entropy flow and vector circulation are
\emph{geometric} rather than phenomenological.

\begin{remark}
The fact that lamphrodynamics sits one degree below the CS example
suggests a natural comparison between topological edge modes and
entropy-flow boundary modes. This will reappear in Chapter~12
(entropic singularities) and Volume~II (black-hole analogs).
\end{remark}



\section{Summary and Forward Pointer}
\label{sec:summary-aksz}

In this chapter we have:
\begin{enumerate}
\item Constructed the AKSZ action for RSVP from the
      $(-1)$--shifted target $\mathrm{Plenum}$.
\item Verified the classical master equation.
\item Identified the BV operator and ghost/antifield expansion.
\item Derived lamphrodynamic PDEs as Euler--Lagrange equations.
\item Described canonical boundary BFV structure.
\item Compared RSVP to the Chern--Simons AKSZ example.
\end{enumerate}

In the next chapter we develop the variational and Hamiltonian
formulation of RSVP in classical degree zero, proving that the
$(-1)$--shifted symplectic geometry induces a well-defined Hamiltonian
system and enabling energy estimates and Noether-type conservation
laws. This prepares the ground for the analytic theory developed in
Chapters~9--11.

